% Hafta 6 — Zamanlama tabanlı sezme
% Derleme: py -3.12 tools/latex/derle.py docs/week-6/assets/h06-08-zamanlama.tex
\documentclass[tikz,border=8pt]{standalone}
\usepackage{cen429-cizim}
\begin{document}
\begin{tikzpicture}
  \node[baslik] (b) at (0,0) {Zamanlama tabanlı sezme};
  \node[altbaslik] at (b.south west) {debugger kodu yavaşlatır};
  \node[kutu=28mm, anchor=north west] (n1) at (0,-2.4) {\kb{t1 = saat()}};
  \node[koyukutu=32mm, right=9mm of n1] (n2) {\kb{Kritik bölge}\\birkaç komut};
  \node[kutu=28mm, right=9mm of n2] (n3) {\kb{t2 = saat()}};
  \node[kotukutu=42mm, right=9mm of n3] (n4)
    {\kb{t2 - t1 > eşik?}\\$\rightarrow$ debugger / kesme noktası};
  \draw[ok] (n1.east) -- (n2.west);
  \draw[ok] (n2.east) -- (n3.west);
  \draw[kotuok] (n3.east) -- (n4.west);
  \node[fit=(n1)(n2)(n3)(n4), inner sep=0pt] (grp) {};
  \node[dip, text width=155mm, anchor=north west] at ($(grp.south west)+(0,-8mm)$)
    {Yanlış pozitif riski: yavaş cihaz, arka plan yükü. Eşiği ölçerek seçin.};
\end{tikzpicture}
\end{document}
